%! Logarithmic Function Context and Data Modeling — Unit 5, Lesson 5.7 — Guided Notes
\documentclass[10pt]{article}
\usepackage{precalculus-article}
\usepackage{precalculus-boxes}
\newcommand{\TallMath}[1]{$\displaystyle #1\rule[-1.4em]{0pt}{3.2em}$}

\begin{document}

\pageheader{Unit 5, Lesson 5.7}{Guided Notes}
\namedateperiod

\vspace{0.1in}

\begin{objectivebox}
\textbf{LO 2.14.A:} Construct a logarithmic function model.

By the end of this lesson I will be able to:
\begin{itemize}[leftmargin=1.5em,itemsep=3pt,topsep=3pt]
  \item Recognize when a context or data set calls for a logarithmic model.
  \item Build a logarithmic model from two input-output pairs using a system of equations.
  \item Apply transformations of $f(x) = a\log_b x$ to fit a model to a context.
  \item Use a logarithmic model to predict values of the dependent variable.
\end{itemize}
\end{objectivebox}

\vspace{0.08in}

\begin{vocabbox}
\termblanklong{Logarithmic model}
\termblanklong{Proportion}
\termblanklong{Real zero}
\termblanklong{Natural logarithm model}
\termblanklong{Logarithmic regression}
\termblanklong{Prediction}
\end{vocabbox}

\vspace{0.08in}

\begin{hookbox}
\textbf{Consider:} A scientist observes that a bacterial population doubles in the 1st hour,
doubles again after 2 more hours, then again after 4 more hours.

\vspace{4pt}
How much time does each additional doubling require? \blank{2cm}, \blank{2cm}, \blank{2cm}

\vspace{4pt}
As the population grows, the time per doubling is \blank{4cm} (increasing / constant / decreasing).

\vspace{4pt}
Why might a logarithmic function --- not an exponential one --- model this situation?
\writeline
\end{hookbox}

\vspace{0.08in}

\begin{notesbox}{When to Use a Logarithmic Model (EK 2.14.A.1)}
A logarithmic model is appropriate when input values change \blank{3cm} over
equal-length output-value intervals.

\vspace{6pt}
\textbf{Contrast:}
\begin{itemize}[leftmargin=1.5em,itemsep=4pt,topsep=2pt]
  \item \textbf{Exponential:} output changes \blank{3cm} over equal input intervals.
  \item \textbf{Logarithmic:} input changes \blank{3cm} over equal output intervals.
\end{itemize}

\vspace{8pt}
\textbf{Example.} Fill in the table and answer the questions.

\vspace{4pt}
\begin{center}
\renewcommand{\arraystretch}{1.8}
\begin{tabular}{|c|c|}
\hline
\rowcolor{sky}
$x$ (area, hectares) & $y$ (number of species) \\
\hline
$2$  & $1$ \\
\hline
$4$  & $2$ \\
\hline
$8$  & $3$ \\
\hline
$16$ & $4$ \\
\hline
\end{tabular}
\end{center}

\vspace{6pt}
As $y$ increases by 1, $x$ \blank{4cm} (describe the pattern). \quad Proportion: \blank{1.5cm}

\vspace{4pt}
This pattern suggests the model: $y = \log_{\blank{1cm}}(x)$
\end{notesbox}

\vspace{0.08in}

\begin{notesbox}{Building a Model from Two Input-Output Pairs (EK 2.14.A.2)}
\textbf{General form:} $f(x) = a\log_b x + k$

\vspace{6pt}
\textbf{Worked Example.} Given: $f(1) = 0$ and $f(e) = 3$. Find the model.

\vspace{4pt}
\textit{Step 1.} Substitute $(1,\,0)$: \quad $0 = a\ln(1) + k = a \cdot \blank{1cm} + k$
\quad so $k = $ \blank{1.5cm}

\vspace{4pt}
\textit{Step 2.} Substitute $(e,\,3)$: \quad $3 = a\ln(e) + 0 = a \cdot \blank{1cm}$
\quad so $a = $ \blank{1.5cm}

\vspace{4pt}
\textit{Model:} $f(x) = $ \blank{4cm}

\vspace{8pt}
\textbf{Practice.} Given: $f(1) = 2$ and $f(9) = 4$. Use base $b = 3$.

\vspace{4pt}
\textit{Step 1.} From $f(1) = 2$: since $\log_3(1) = $ \blank{1cm}, we get $k = $ \blank{1.5cm}

\vspace{4pt}
\textit{Step 2.} From $f(9) = 4$: $4 = a\log_3(9) + k = a \cdot \blank{1cm} + \blank{1cm}$
\quad so $a = $ \blank{1.5cm}

\vspace{4pt}
\textit{Model:} $f(x) = $ \blank{4cm}
\end{notesbox}

\vspace{0.08in}

\begin{notesbox}{Transformations of a Logarithmic Model (EK 2.14.A.3)}
General form with transformations: $f(x) = a\log_b(x - h) + k$

\vspace{4pt}
\begin{itemize}[leftmargin=1.5em,itemsep=4pt,topsep=2pt]
  \item $a$: vertical \blank{3cm} (stretch if $|a|>1$, compression if $|a|<1$)
  \item $h$: horizontal shift \blank{1cm} units to the \blank{2cm}
  \item $k$: vertical shift \blank{1cm} units \blank{2cm}
\end{itemize}

\vspace{6pt}
\textbf{Context example.} A sound-level function has output 0 dB at intensity $x = 1$
and increases by 10 dB for every tenfold increase in intensity.

\vspace{4pt}
The output is 0 when $x = 1$, so \blank{3cm} (which parameter?) equals \blank{1cm}.

\vspace{4pt}
A tenfold increase in $x$ adds 10 to the output, so $a = $ \blank{1.5cm} and $b = $ \blank{1.5cm}.

\vspace{4pt}
Model: $S(x) = $ \blank{5cm}
\end{notesbox}

\vspace{0.08in}

\begin{notesbox}{Predicting with a Logarithmic Model (EK 2.14.A.6)}
Given $f(x) = 2\log_3 x$:

\vspace{6pt}
\textbf{Predict $f(81)$:} \quad $f(81) = 2\log_3 81 = 2 \cdot \blank{1.5cm} = $ \blank{2cm}

\vspace{6pt}
\textbf{Solve for $x$ when $f(x) = 5$:}

\vspace{4pt}
$2\log_3 x = 5 \;\Rightarrow\; \log_3 x = \blank{1.5cm} \;\Rightarrow\; x = 3^{\blank{1.2cm}} = $ \blank{3cm}

\vspace{8pt}
\textbf{Natural log models (EK 2.14.A.5).} Recall $\ln x = \log_e x$.

\vspace{4pt}
Useful identities: \quad $\ln(e^n) = $ \blank{1.5cm} \quad and \quad $e^{\ln x} = $ \blank{1.5cm}

\vspace{4pt}
If $f(x) = 3\ln x$, evaluate $f(e^4) = 3\ln(e^4) = 3 \cdot \blank{1.5cm} = $ \blank{2cm}

\vspace{4pt}
Solve $3\ln x = 9$: \quad $\ln x = $ \blank{1.5cm} \quad $\Rightarrow$ \quad $x = e^{\blank{1.2cm}} = $ \blank{2cm}
\end{notesbox}

\end{document}
